\documentclass{article}
\usepackage{graphicx}
\usepackage{amsmath}
\usepackage{amsfonts}
\usepackage{bookmark}
\usepackage{booktabs}
\usepackage{exercises}
\graphicspath{ {./images/test} }
\usepackage{indentfirst}
\usepackage{amsthm}
\usepackage{yhmath}
\usepackage{placeins}
\newcommand*{\oldneg}{\mathord{\sim}}
\newcommand\showdiv[1]{\overline{\smash{)}#1}}
\newcommand\Mydiv[1]{%
\smash{\raise.3ex\hbox{$\big)$}}\mkern-8mu
        \overline{\enspace\strut#1}}
\setlength{\tabcolsep}{5pt}
\theoremstyle{definition}
\newtheorem{definition}{Definition}

\setlength{\parindent}{.5in}
\setlength{\parskip}{12pt}
\renewcommand{\arraystretch}{2}
\newcommand{\?}{\stackrel{?}{=}}

\title{Mixed Review of Problem Solving (pg. 216)}
\begin{document}
\maketitle

\section*{1}
(a): \[ R=1.4C \]
(b): \[ R=1.4C \] \[ \frac{R}{1.4}=C \]
(b): \[ \frac{60}{1.4}=C \] \[ C\approx 42.8 \]

\section*{2}
\[ \sqrt{3x-5}=4 \]
\[ 3x-5=16 \]
\[ x=7 \]
To use the graph to solve this equation, we can see what value the graph has at $x=7$. Looking at the graph, the value is $y=4$, 
confirming that our initial solution was correct.

\section*{3}
\[ \sqrt[3]{x}=\sqrt[3]{-5} \]

\section*{4}
(a): \[ k=7.463d\textrm{ (Rewrite given equation)} \] \[ d=\frac{k}{7.463} \]
(b): \[ d=\frac{25}{7.463}\approx \$ 3.350 \]
(c): The inverse functions tells us hom many U.S. dollars you recieve for a given amount $k$ of kronor.

\section*{5}
We know that $y=\sqrt{x}$ passes through $0,0$ and $1,1$. For a given translation of a radical equation of the form:
\[ a\sqrt{x-h}+k \] the graph of this function is shifted $h$ units to the right and $k$ units up from the graph of $y=\sqrt{x}$. 
So, we can take the function $\sqrt{x+14}$, which would shift $\sqrt{x}$ over to the left by $14$ units, meaning it passes through 
$(-14,0)$ and $(-13,1)$.

\section*{6}
First function: $y=\sqrt{x-4}$ \newline
Second function: $y=\sqrt{x}-2$

\section*{7}
(a): \[ \frac{v^2}{64}=h \]
(b): Suppose that we are given an initial velocity of 32 feet per second, and we wish to find the maximum height. We can deduce that:
\[ h=\frac{32^2}{64}=\frac{1024}{64}=\frac{2^{10}}{2^6}=2^4=16 \] feet.

\section*{8}
We are given that $d=100$, meaning:
\[ 100=\sqrt{2500+h^2} \]
\[ 10000=2500+h^2 \]
\[ 7500=h^2 \]
\[ h\approx 86.6 \] feet.

\section*{9}
\[ A(r(t))=A(6t)=\pi(6t)^2=36t^2\pi \]
meaning \[ A(r(2))=36(2)^2\pi=144\pi \] $A(r(2))$ represents the area of the circle that is formed from the radius of the circle that is 
formed two seconds after the pebble hits the water.

\end{document}
